\documentclass[12pt]{book}
\usepackage[utf8]{inputenc}
\usepackage{xskak}
\usepackage{amsfonts}
\usepackage{amsmath}
\usepackage{amsthm}
\usepackage{chessboard}
\usepackage[margin=1in]{geometry}
\usepackage{longtable}
\usepackage{booktabs}
\usepackage{hyperref}
\usepackage{float}

% Custom colors for move quality
\definecolor{brilliantcolor}{RGB}{0, 150, 150}
\definecolor{excellentcolor}{RGB}{0, 128, 0}
\definecolor{goodcolor}{RGB}{64, 160, 64}
\definecolor{inaccuracycolor}{RGB}{200, 180, 0}
\definecolor{mistakecolor}{RGB}{220, 120, 0}
\definecolor{blundercolor}{RGB}{200, 0, 0}

\def\FEN{\tt{FEN}}
\def\qqq{\mathbb{Q}}
\def\rrr{\mathbb{R}}
\def\zzz{\mathbb{Z}}

\newtheorem{example}{Example}[section]

\title{Chess Game Analyzer,\\
an introduction}
\author{David Joyner}
\date{2026-02-02}

\begin{document}
\maketitle

\tableofcontents

\newpage


\chapter*{Preface}

This is a guide of how to use a program I wrote to
analyze chess games you input (as a pgn file on your computer)
called the {\it chess game analyzer}, or {\it CGA} \cite{CGA26}.
The program uses other chess programs to output a
math-dense LaTeX file, called a
{\it Quantitative-Analytic Report on Chess} (or QARC),
where the results from those other programs are collected
in a human-readable format. Full disclosure: For creating CGA,
I did get a lot of help from AI: {\tt gemini} (Google) for
chess-theoretic questions, {\tt chatGPT} (openAI) for outline ideas,
and {\tt claude} (Anthropic) for Python questions.

The CGA is a Python module providing
an advanced chess analysis pipeline that goes beyond piece
evaluations. It combines {\tt Stockfish} engine analysis with
``home-made'' positional metrics to generate quantitative
LaTeX reports on chess games, machine-readable data, and
win-prediction metrics.
In addition to generating educational content, it collects
useful game data into tables, and computes
``Fireteam Index'' (or FTI, defined below) signatures.
As we will see below, the FTI is a family of metrics inspired by
the medieval board game rithmomachia \cite{Ri25} to measure
sustained positional dominance.

In fact, all this, the CGA module and related
programs\footnote{For instance, CGA has an add-on that creates
animated gifs of games, but such add-ons will be discussed separately.},
was inspired by the ``territorial'' approach that proper
victories take in rithmomachia. Can such an approach be helpful in
chess analysis? Maybe more refinements to CGA are needed before I can answer
that, but that question is the main motivation for this work.

Another goal was to create ``Prediction Algorithms'' for chess games
in the hope of predicting a win during live play.
Of course, that is impossible in general, but the question remains:
what can be done,
provided we allow some (not too many) inaccuracies?
To address this, CGA implements three algorithms that use the FTI to
predict game outcomes:

\begin{itemize}
\item
  {\tt Per-Ply}.
 \begin{itemize}
 \item
  Examines FTI at each ply (=half-move) after ply 16.
 \item
   Predicts {\it win} if FTI $> \epsilon$ (default $\epsilon = 0.2$)
   for 10 consecutive plies.
 \item
  Added plus: Captures sudden, decisive advantages.
 \end{itemize}
\item
  {\tt Windowed}
 \begin{itemize}
 \item
 This is a ``smoothed'' version of {\tt Ply}.
 \item
  Computes a 10-ply rolling average of FTI, FTI$_{10}$ say.
 \item
   Applies the same streak detection to smoothed values:
   Predicts {\it win} if FTI$_{10} > \epsilon$ 
   for 10 consecutive plies.
 \item
  Filters out tactical noise; captures sustained trends.
 \end{itemize}
\item
  {\tt Balance} 
 \begin{itemize}
 \item
 Counts plies where FTI $> \epsilon$ versus FTI $<- \epsilon$
 \item
 Measures volatility (sign changes in the difference between these counts)
 \item
   Computes confidence = $|B_n|\cdot (1 - V/V_{\rm max})$
\item   
  Predicts {\it draw} if confidence $<0.15$; otherwise {\it win/loss}
  (that is, not a draw).
\item
  Best at detecting drawn games and balanced struggles
 \end{itemize}
 \end{itemize}

The example games are well-known and well-analyzed, so you can
compare CGA analysis here with the GM analysis published already
in books, or blogs or on youtube.
The first one is a World Correspondence Chess game played in 1965,
Estrin-Berliner, \S \ref{sec:estrin-berliner_1965}.
The second is Game 6 of the World Chess Championship match
Carlsen-Caruana in London, Nov. 2018, \S \ref{sec:carlsen-caruana_game6-2018}.

You don't need to know how to use {\tt Stockfish} or {\tt Python}
in order to run the CGA programs. However, you do need to know how to
install them on your computer if they aren't already.
Installation help and the syntax for using CGA is presented
in detail in the appendix,
ch. \ref{appendix:cga-readme}. Don't worry about the syntax or
installing software yet. First, let's learn a little more
about the CGA module.


\chapter{Introduction}

As mentioned,
detailed reference for the CGA module is given in an appendix,
ch. \ref{appendix:cga-readme}.
However, the basic idea is that the CGA module takes a pgn file as input.
The output consists of a latex report which contains the following
information about the game (or games) in the pgn file.

\begin{itemize}
\item  
Move-by-move positional analysis with Space, Mobility, King Safety, and
Threats metrics
\item  
{\tt Stockfish} suggestions of playable alternative moves within 50 centipawns
of the best 
\item  
  Game character classification (balanced/tense/tactical/chaotic,
  one-sided/seesaw)
\item  
Brilliant sacrifice detection with sound/speculative classification
\item  
Five different Fireteam Index metrics capturing different playing styles
\item
Other metrics, such as Balance and Volatility, measure liklihood of a draw.
\item  
Three prediction algorithms (Per-Ply, Windowed, Balance) for win or draw predictions.
\item  
Publication-quality LaTeX output with chess diagrams, plots, and methodology
\item  
Book mode for analyzing multi-game PGN files with aggregate statistics
\item  
Raw data export for downstream computational analysis

\end{itemize}

These will be explained, with examples, below.

\section{Computer chess for humans}

When a human plays chess in a given position, they use
both logic (to trace through potential moves and the
ramifications of their opponent's responses) and
experience (to restrict their analysis
to a small number of moves) in their effort to win.
The human player makes a move based on rigorous reasoning and their
previous experience with games similar to the one they are
currently playing.
They pick one that, they feel, gives them the best position no
matter what their opponent does in response. Of course,
a player will not want to lose any advantage they feel they
have in the position, whether material, or space, or mobility, or king
safety, or potential threats they can mount.
At the same time, they want their move to also
restrict their opponent's advantages in the position.

A computer plays chess using a different process.
The computer player also uses logic to sift through legal moves
for the best one. The difference is that the computer can, with the help of
a chess engine like {\tt Stockfish}, evaluate positions very rapidly
and they pick the move that both maximizes their own position
evaluation and minimizes the evaluations of each response by their opponent.
In other words, to a computer, playing chess is analogous to
solving a math problem. Can humans learn a lesson or two from how
the computer evaluates positions? Let's look more closely at this
algorithmic perspective the computer player takes.

A real-valued function $m:\FEN\to\rrr$
that takes in a position of a chess game and
returns a positive number is called a {\it chess metric}.
Here, $\FEN$ simply denotes the collection of
all possible FEN positions\footnote{FEN stands for
Forsyth–Edwards Notation, a standard notation for describing a
board position, typically represented as an alphanumeric list or string,
see \cite{FEN}. Forsyth was a Scottish journalist from the
1800s while Edwards was an American computer scientist active in
computer chess from the late 1900s to early 2000s.}
and $\rrr$ is the standard notation for the set of real numbers.

One easy metric all players use is the {\it piece value metric}
for White and for Black: $PV_W({\rm fen})$ denotes the total value
of the white pieces in the position ${\rm fen}$,
and $PV_B({\rm fen})$  using the standard values for the Black pieces
in the give position.
Most players use the usual values\footnote{For more background
on the different ways to evaluate chess pieces, and there
are quite a lot of them, see 
the wikipedia article \cite{PV} and the references there
(including references to Berliner's book \cite{Be99}).} for the
pieces:

\begin{center}
\begin{tabular}{lc}
piece & value \\ \hline
king & $\infty$ \\
queen & $9$ \\
rook & $5$ \\
bishop & $3$ \\
knight & $3$ \\
pawn & $1$ \\ \hline
\end{tabular}
\end{center}

\noindent
The reader can check that, in the game's starting position,
we have

\[
P_W({\rm start}) = P_B({\rm start}) = 39.
\]

Another example is a chess metric we'll denote
by {\tt eval}, the evaluation metric from
the chess engine {\tt Stockfish}.
This positional metric measures advantages (including piece values) in
{\it centipawns} (cp) —-- hundredths of a pawn.
An evaluation of $+100$ (that is, ${\tt eval}({\rm fen})>100$)
means White is ahead in the given \FEN position
by the positional equivalent of roughly one pawn. An evaluation of
$-250$ means Black has about a $2.5$ pawn advantage.

\begin{center}
\begin{tabular}{|c|c|}
Evaluation & Meaning \\ \hline
 $0$ to $\pm 50$ & Roughly equal \\
 $\pm 50$ to $\pm 150$ & Slight advantage \\
 $\pm 150$ to $\pm 300$ & Clear advantage \\
 $\pm 300$ to $\pm 500$ & Winning advantage \\
 $\pm 500+$ & Decisive advantage \\
$\pm 10000$ & Checkmate is forced \\
\end{tabular}
\end{center}

The modern {\tt Stockfish} versions primarily uses NNUE (neural network)
evaluation to compute {\tt eval}. Unfortunately, that metric is a bit of
a black box. It appears to be impossible, or at least very impractical, to
directly extract from the modern {\tt Stockfish} other,
more ``territorial'' metrics such as Space.

Let's dive into that topic next.


\section{The five core metrics}

The {\tt Stockfish} evaluation {\tt eval} tells the story, metrics explain why.

As mentioned,
the modern {\tt Stockfish} {\tt eval} is a bit of a black box.

Given a chess position, who do we think of it from a territorial
perspective? These natural questions arise:

\begin{itemize}
\item
How to we compute the ``space'' controlled by a given side?

\item
How to we compute the ``mobility'' of pieces in the position?

\item
How do we compute the safety of the king?
This is, once it's defined, a defensive metric.

\item
How do we compute the ``threats'' one side has against their opponent?
Conversely, once it's defined, this is an offensive metric.
\end{itemize}

The four metrics
Space/Mobility/King Safety/Threats use the Python module
{\tt python-chess} to parse the position and then Python to directly
compute the metrics based on reasonable
formulas\footnote{For details on {\tt python-chess} and the
formulas for the metrics (Space, Mobility, and so on),
please see the following subsections or the
appendix, chapter \ref{appendix:cga-readme}.}.
These formulas are designed to address these questions.

\subsection{Data model and metric plots}

Raw positional data fields are printied to the bottom of
the latex QARC file (but does not appear in the
compiled pdf). CGA uses this data to plot various chess metrics.

The plots of the metrics are either Per-ply or Windowed.

\begin{itemize}
\item
Per-ply: Looks for 10 consecutive positions where one side dominates
(FTI $> 0.20$). 
\item
Windowed/Smoothed: Uses a 10-ply rolling average to filter out tactical noise.
\end{itemize}

\subsection{Space}

``How much `room' do your pawns give your pieces to move in the center?''

Space evaluation counts squares in the center (files c-f, ranks 2-4 for White, 5-7 for Black) that are:
\begin{enumerate}
\item Attacked by at least one friendly pawn, {\it and}
\item Not occupied by any enemy piece
\end{enumerate}

The computation is weighted by the number of pieces behind the pawn
chain --- this rewards positions where you control territory {\it and}
have pieces active in that space.
The Space term is scaled by total piece count; it matters more in
closed positions with many pieces.

Space is a middlegame metric --- it becomes near-meaningless in endgames when
most pieces are traded.

\subsection{Mobility}

``How many useful squares can your pieces actually land on? (More squares = more power)''


Mobility counts legal moves for each piece type, with different weights:
\begin{itemize}
\item \textbf{Knights}: Safe squares attacked (excluding squares defended by enemy pawns)
\item \textbf{Bishops}: Safe squares on diagonals (bonus for long, unobstructed diagonals)
\item \textbf{Rooks}: Squares on files/ranks (bonus for open files, 7th rank penetration)
\item \textbf{Queens}: Combination of bishop and rook mobility patterns
\end{itemize}

``Safe squares" exclude those defended by enemy pawns or where exchanges would lose material.
Mobility in the opponent's territory is valued higher than mobility in your own.

Mobility tracks piece activity. Watch it climb during development and converge as
long-range pieces (bishops, rooks, queens) are trade off.
  
\subsection{King Safety}

`How well is your King shielded by pawns, and how many enemy `attack units' are nearby?''

King safety combines several factors:
\begin{itemize}
\item \textbf{Pawn shield}: Pawns on 2nd/3rd rank directly in front of the castled king
\item \textbf{King geometry}: Distance of enemy pieces to your king.
  
\item \textbf{Attack units}: Count of enemy pieces attacking squares adjacent to your king
\item \textbf{Safe checks}: Number of safe checking squares available to the opponent
\end{itemize}

King Safety can be positive and negative. Negative values in the middlegame mean danger. Negative values in the endgame can mean active, fighting kings.

\subsection{Threats}

``Are you hanging pieces, or are your small pieces attacking the opponent's big
pieces?''

Threat evaluation considers:
\begin{itemize}
\item Hanging pieces (attacked but not adequately defended)
\item Safe pawn advances that create threats
\item Minor pieces attacking major pieces
\item Weak squares in the pawn structure
\end{itemize}

Threats measures potential piece captures. High threat values show tension;
when they drop sharply after exchanges, material or positional concessions
have been made.


\begin{example}
For example, the $-319$cp blunder at move 15 in \S \ref{sec:estrin-berliner-15Be2}
below shows what happened but the Mobility and Threats changes in moves 8-14
show why the position was ripe for sharp play form both sides.
\end{example}

\subsection{The Fireteam Index signature}

This is a linear combination of those four metrics above.

The Fireteam Index is a composite metric inspired by the ancient mathematical
board game {rithmomachia}, where a ``progression victory'' is achieved when a
coordinated group of pieces (a ``fireteam'') establishes control deep in enemy
territory \cite{Ri25}. Similarly, in chess, sustained positional dominance
across multiple factors often precedes victory.

The FTI weights can be thought of as \emph{signatures} that capture different
playing styles. By adjusting which positional factors receive emphasis,
the FTI can be tuned to better predict outcomes for players with different strategic preferences.



\subsection{Balance}

Balance is the key draw-detection metric. It combines:

\begin{itemize}
\item
Normalized Balance ($B_n$): What fraction of positions favored each side
\item
Volatility (V): How often the advantage changed hands
\item
Confidence (C): $C=B_n\cdot (1 - V/V_{\rm max})$, where $V_{\rm max}$ is
the maximum volatility for that game.
\end{itemize}

The idea is that
if the positional metrics are roughly equal for most of the game
then neither side dominated the Space/Mobility/Threats picture consistently.
Similarly, high volatility means the advantage kept changing hands
during the course of the game. The {\it Balance algorithm} will predict
the game ends in a draw if $|B_n|<0.2$ or $V \geq 10$.
This data is part of the QARC output.

\chapter{Examples}

The first example game, ironically, was won by
Hans Berliner, a highly respected computer pioneer and computer science
researcher.

\section{Estrin-Berliner, 1965}
\label{sec:estrin-berliner_1965}
        
This is a popular game for commentary and analysis.
The point here is not to try to provide a summary of the
analysis of others but to report on what CGA says about
these games.

The QARC file on the game Estrin vs. Berliner (1965) describes
it as a ``chaotic seesaw battle,''
where the advantage changed hands multiple times \cite{Be65}.
 We'll walk you through this game using the provided metrics, you can use them
to explain the dynamics of the positional components (Space, Mobility, King Safety, and Threats).

White, played by the Russian chess player Yakov Estrin, took joint first place in the
USSR Correspondence Championship
a few years earlier, in 1962.
Estrin is best known for being the seventh ICCF World Champion a few
years later, in 1972–1976.

Black, played by American Hans Berliner, was also a well-known chess player.
Here is a quotation from Dunne's rememberance article for
Has Berliner written for {\bf Chess Life} \cite{Du17}:

\begin{quotation}
The finals of the Fifth World Correspondence Chess Championship
produced a score of $14-2$, $12$ wins and four draws, a three point margin
over the second place finisher, a spread unmatched in World Championship
play. In that final, Hans also played what is widely regarded as
the greatest correspondence game ever played, and listed in Andy Soltis’
book {\bf The 100 Best Chess Games of the 20th Century} as the top game of
that tumultuous century. It is a game that has resisted major
improvement on Black’s part for over half a century.
\end{quotation}

In fact, that last statement of Dunne's is, thanks to progress in
chess engines such as {\tt Stockfish}, no longer completely true.
However, we aren't going down that rabbit hole in this guide.

\subsection{Opening phase (Moves 1–5): The Two Knights Defense}

With that background out of the way, let's see what the CGA
had to say about this game.

\vskip .1in
\noindent
{\bf 
1.e4 e5 2.Nf3 Nc6 3.Bc4 Nf6 4.Ng5 d5 5.exd5
}

\vskip .1in
\noindent
This is the famous Two Knights Defense with the Fried Liver Attack.
Estrin plays 4.Ng5,
immediately creating tension. Let's see how the metrics capture this.

Notice how Space remains roughly equal at 0.40 for both sides through move 9.
Space measures control of central squares backed by your pawns —- essentially,
how much room your pieces have to maneuver.
After

\vskip .1in
\noindent
{\bf 
5.exd5
}

\vskip .1in
\noindent
White captures a pawn but Black will strike back. Watch what happens after
5...b5 --— the Ulvestad Variation.
The usual move in this position is 5. ... Na5.

%\begin{center}
\begin{figure}[h!]
\begin{minipage}{\textwidth}
\begin{center}
%\begin{tabular}{c}
%\hspace{4.0 cm}
  %\vspace{1.0 cm}
\chessboard[setfen=r1bqkb1r/p1p2ppp/2n2n2/1p1Pp1N1/2B5/8/PPPP1PPP/RNBQK2R w KQkq - 0 6]  
%\includegraphics[height=7cm,width=7cm]{berliner-game7-position1-after5-b5.png}
\end{center}
%\\
\begin{center}
%\hspace{4.0 cm}
\includegraphics[height=7cm,width=12cm]{berliner-game7-threats-metric.png}
%\end{tabular}
\end{center}
\end{minipage}
\caption{(top) Position after 5 ... b5 in Estrin-Berliner (1965). (bottom) The Threats metric for that game. }
\label{fig:estrin-berliner-position1}
\end{figure}
%\end{center}

Data at ply 10/move 5:
%
%\begin{itemize}
%\item  
%Space: W 0.40, B 0.30
%\item  
%Mobility: W 4.40, B 5.10
%\item  
%King Safety: W 1.00, B -0.10
%\item  
%Threats: W 10.80, B 6.80
%\item  
Eval: +0.45
%\end{itemize}

Look at the Threats metric -— White's threats spike to 10.80 because the b5-pawn is hanging.
But if you compare Black's Mobility to White's, you see
Berliner is accepting material deficit for piece activity.


\subsection{The critical opening battle (Moves 6–15)}
\label{sec:estrin-berliner-15Be2}

\vskip .1in
\noindent
{\bf
6. Bf1 Nd4 7. c3! Nxd5! 8. Ne4?! Qh4?
}

\vskip .1in
\noindent
From the game's annotation by {\tt Stockfish}, both 8 ... Nc6 and 8 ... Ne6 are preferred.

\vskip .1in
\noindent
{\bf
9. Ng3! Bg4?
}

\vskip .1in
\noindent
Again, {\tt Stockfish} recommends 9 ... Nc6.

\begin{center}
\begin{figure}[h!]
\begin{minipage}{\textwidth}
\begin{center}
%\begin{tabular}{c}
%\hspace{4.0 cm}
  %\vspace{1.0 cm}
\chessboard[setfen=r3kb1r/p1p2ppp/8/1p1np3/3n2bq/2P3N1/PP1P1PPP/RNBQKB1R w KQkq - 4 10]  
%\includegraphics[height=7cm,width=7cm]{berliner-game7-position2-after9-Bg4.png}
\end{center}
%\\
\begin{center}
%\hspace{4.0 cm}
\includegraphics[height=7cm,width=12cm]{berliner-game7-mobility-metric.png}
%\end{tabular}
\end{center}
\end{minipage}
\caption{(top) Position after 9 ... Bg4 in Estrin-Berliner (1965). (bottom) The Mobility metric for that game. }
\label{fig:estrin-berliner-position2}
\end{figure}
\end{center}

Data at ply 18/move 9:

%\begin{itemize}
%\item  
%Space: W 0.30, B 0.30
%\item  
%Mobility: W 3.90, B 5.00
%\item  
%King Safety: W -0.40, B 0.50
%\item  
%Threats: W 3.50, B 9.50
%\item  
Eval: +1.55
%\end{itemize}

The game goes on (see \cite{Be65}):

\vskip .1in
\noindent
{\bf
10. f3! e4?
}

\vskip .1in
\noindent
Black loses 67cp. Consider instead: Bd7, Be6, Bd6.

\vskip .1in
\noindent
{\bf
11. cxd4 Bd6! 12. Bxb5+ Kd8 13. O-O exf3 14. Rxf3! Rb8 
}

\vskip .1in
\noindent
This sets up a crutial position.


\subsection{The blunder}

With threats from both sides, the next move is important.

\vskip .1in
\noindent
{\bf
15. Be2??
}

\vskip .1in
\noindent
{\it Looks} good but, in fact, some commentators say
this natural-looking move is the move that lost the game.
Stockfish recommends 15 Bf1 or else 15 a4.


\begin{center}
\begin{figure}[h!]
\begin{minipage}{\textwidth}
\begin{center}
%\begin{tabular}{c}
%\hspace{4.0 cm}
%\vspace{1.0 cm}
\chessboard[setfen=1r1k3r/p1p2ppp/3b4/3n4/3P2bq/5RN1/PP1PB1PP/RNBQ2K1 b - - 2 15]
%\includegraphics[height=7cm,width=7cm]{berliner-game7-position3-after15Be2.png}
%\end{center}
%\\
%\begin{center}
%\hspace{4.0 cm}
%\includegraphics[height=7cm,width=12cm]{berliner-game7-mobility-metric.png}
%\end{tabular}
\end{center}
\end{minipage}
\caption{Position after 15 Be2 in Estrin-Berliner (1965).}
\label{fig:estrin-berliner-position3}
\end{figure}
\end{center}

Data at ply 31:

%\begin{itemize}
%\item  
%Space: W .30, B 0.30
%\item  
%Mobility: W 5.00, B 4.60
%\item  
%King Safety: W 0.00, B 0.10
%\item  
%Threats: W 3.80, B 4.70
%\item  
Eval: -1.26
%\end{itemize}

\subsection{The tactical middlegame (Moves 15–28)}

The game continues:

\vskip .1in
\noindent
{\bf
  16. Bxf3! Qxd4+?
}
\vskip .1in

\noindent
{\tt Stockfish} preferred instead 16... Re8.

\vskip .1in
\noindent
{\bf
17. Kh1! Bxg3! 18. hxg3! Rb6 19. d3! Ne3! 20. Bxe3! Qxe3! 21. Bg4 h5 22. Bh3! g5
23. Nd2! g4! 24. Nc4! Qxg3! 25. Nxb6! gxh3! 26. Qf3 hxg2+ 27. Qxg2! Qxg2+
}

\vskip .1in
\noindent
The queens are trading off now.

\subsection{The endgame, initial phase (Moves 28–34)}

Without long-range pieces, the Mobility metric drops.
We enter a new phase of the endgame.

\vskip .1in
\noindent
{\bf
28. Kxg2 cxb6! 29. Rf1 Ke7! 30. Re1+! Kd6 31. Rf1! Rc8?
}

\vskip .1in
\noindent
{\tt Stockfish} recommends Ke6, Ke7, or Rh7 instead.

\begin{center}
\begin{figure}[h!]
\begin{minipage}{\textwidth}
\begin{center}
%\begin{tabular}{c}
%\hspace{4.0 cm}
%\vspace{1.0 cm}
\chessboard[setfen=2r5/p4p2/1p1k4/7p/8/3P4/PP4K1/5R2 w - - 6 32]
%\includegraphics[height=7cm,width=7cm]{berliner-game7-position4-after31-Rc8.png}
\end{center}
%\\
\begin{center}
%\hspace{4.0 cm}
\includegraphics[height=6cm,width=12cm]{berliner-game7-space-metric.png}
%\end{tabular}
\end{center}
\end{minipage}
\caption{(top) Position after 31 ... Rc8? in Estrin-Berliner (1965). (bottom) The Space metric for that game. }
\label{fig:estrin-berliner-position4}
\end{figure}
\end{center}


Data at ply 62/move 31:

%\begin{itemize}
%\item  
%Space: W 0.30, B 0.20
%\item  
%Mobility: W 1.80, B 1.80
%\item  
%King Safety: W -0.20, B -1.20
%\item  
%Threats: W 4.50, B 2.10
%\item  
Eval: -1.21
%\end{itemize}

\vskip .1in
\noindent
{\bf
32. Rxf7?! Rc7! 33. Rf2?
}

\vskip .1in
\noindent
{\tt Stockfish} suggests to consider instead the move Rf6+.


\vskip .1in
\noindent
{\bf
33...Ke5 34. a4??
}

\vskip .1in
\noindent
This is a more serious error, but not of the same magniture as 15 Be2??.
With 34. a4??, {\tt eval} drops from $-2.24$ to $-4.01$.
Here, {\it Stockfish} recommends instead Rf8.

\vskip .1in
\noindent
{\bf
34...Kd4! 35. a5?! Kxd3
}


\subsection{The final phase (Moves 35–42)}

The final moves of the game.

\vskip .1in
\noindent
{\bf
36. Rf3+ Kc2?! 37. b4 b5! 38. a6?! Rc4?! 39. Rf7 Rxb4 40. Rb7?! Rg4+! 41. Kf3?
}

\vskip .1in
\noindent
{\tt Stockfish} suggests instead Kh3 or Kf3.

\vskip .1in
\noindent
{\bf 
41...b4?! 42. Rxa7! b3! 0-1
}

\vskip .1in


\subsection{FTI summary}

The various types of FTIs are designed to characterize the type of play.
Recall that the FTI metric is a linear combination of the four positional metrics,
Space (S), Mobility (M), King Safety (KS), and Threats (T):

\[
FTI = w_SS + w_M M + w_{KS}KS + w_T T.
\]
Think of the FTI as a coordinated pressure gauge.
Using it to predict a game's outcome requires a player to hold
positional dominance for 10 consecutive 'plies' (half-moves)
to predict a win. For the Estrin-Berliner game, predicting the outcome
(win/loss/draw) using Fireteam Indices is a struggle because the game was so volatile.


In this case, CGA was run over all the games Berliner played for the
1965 World Correspondence Chess championship. The CGA program computes the
weights $(w_S, w_M, w_{KS}, w_T)$ such that the FTI with those
``signature'' weights, the plot predicts the outcomes more successfully than
any other choices. The outcome was

\[
(w_S, w_M, w_{KS}, w_T) = ( 0.319,  0.675,  0.003,  0.003 ),
\]
with $R^2 = -12.525$ (indicating a poor fit).

Either by looking at the plots for FTI3 and FTI5 or that signature FTI plot,
we see CGA suggests Berliner has a solid, strategic style.


\begin{center}
\begin{figure}[H]
\begin{minipage}{\textwidth}
\begin{center}
%\begin{tabular}{c}
%\hspace{4.0 cm}
%\vspace{1.0 cm}
\includegraphics[height=7cm,width=12cm]{estrin-berliner-1965-FTI1.png}
\end{center}
%\\
\begin{center}
%\hspace{4.0 cm}
\includegraphics[height=7cm,width=12cm]{estrin-berliner-1965-FTI2.png}
%\end{tabular}
\end{center}
\end{minipage}
\caption{Both plots are for Estrin-Berliner, 1965. (top) Fireteam index, FTI1 (harmonious). (bottom) Fireteam index, FTI2 (tactical). The second plot never exceeds the confidence threshold of $\pm 0.2$.}
\label{fig:estrin-berliner-1965-FTI1+2}
\end{figure}
\end{center}

\begin{center}
\begin{figure}[H]
\begin{minipage}{\textwidth}
\begin{center}
%\begin{tabular}{c}
%\hspace{4.0 cm}
%\vspace{1.0 cm}
\includegraphics[height=7cm,width=12cm]{estrin-berliner-1965-FTI3.png}
\end{center}
%\\
\begin{center}
%\hspace{4.0 cm}
\includegraphics[height=7cm,width=12cm]{estrin-berliner-1965-FTI4.png}
%\end{tabular}
\end{center}
%\\
\begin{center}
%\hspace{4.0 cm}
\includegraphics[height=7cm,width=12cm]{estrin-berliner-1965-FTI5.png}
%\end{tabular}
\end{center}
\end{minipage}
\caption{These plots are also for Estrin-Berliner, 1965. (top) Fireteam index, FTI3 (strategic). (middle) Fireteam index, FTI4 (dynamic). (bottom) Fireteam index, FTI5 (solid).}
\label{fig:estrin-berliner-1965-FTI3+4+5}
\end{figure}
\end{center}

Of all the FTI plots in Figures \ref{fig:estrin-berliner-1965-FTI1+2}
and \ref{fig:estrin-berliner-1965-FTI3+4+5}, the last one is the one that
most strongly suggests a win for Black. FTI1 and FTI3 also do suggest a
Black win, but with less confidence.


\section{Carlsen-Caruana, Game 6, 2018}
\label{sec:carlsen-caruana_game6-2018}

At this point in the match, the players were tied
2.5-2.5, with 5 draws in a row. This game was another draw but exciting none-the-less.
In fact, on move 67 it appears Caruana missed a win.

The QARC file on the $12$ games in the Carlsen-Caruana match, London
(November 2018) is available online as a pdf \cite{CC18}.
We'll walk you through this game using the provided metrics,
you can use them
to explain the dynamics of the positional components
(Space, Mobility, King Safety, and Threats).

Of course, both players are well-known.
In 2018, the 27-year-old Norwegian was the undisputed king of chess,
having held the World Championship title since 2013 and the world
No. 1 ranking since 2011.
The 26-year-old American-Italian grandmaster, in 2018, became
the first American to challenge for the undisputed world chess title since
Bobby Fischer in 1972.
At the time of their match, their ELO ratings were nearly the same.


\subsection{The early opening moves (1-4)}

Look at the plot for {\tt Stockfish}'s {\tt eval} function (see Figure \ref{fig:carlsen-caruana-position1}).


\begin{center}
\begin{figure}[h!]
\begin{minipage}{\textwidth}
\begin{center}
%\begin{tabular}{c}
%\hspace{4.0 cm}
%\vspace{1.0 cm}
%\includegraphics[height=7cm,width=7cm]{berliner-game7-position4-after31-Rc8.png}
\chessboard[setfen=rnbqkb1r/ppp2ppp/3p1n2/8/4P3/3N4/PPPP1PPP/RNBQKB1R b KQkq - 1 4]
\end{center}
%\\
\begin{center}
%\hspace{4.0 cm}
\includegraphics[height=7cm,width=12cm]{carlsen-caruana-game6-eval-metric.png}
%\end{tabular}
\end{center}
\end{minipage}
\caption{(top) Position after 4 Nd3?! in Carlsen-Caruana, Game 6 (2018). (bottom) The {\tt eval} metric for that game. }
\label{fig:carlsen-caruana-position1}
\end{figure}
\end{center}

Notice how the line hovers around zero for the first 10-15 moves,
then drifts into negative territory,
favoring Caruana. The game reaches its most critical point around
moves 67-68 where we see a
dramatic spike to about $-2.5$ pawns, a huge advantage for Black.
Unfortunately for Caruana, that was followed by an equally dramatic
swing back toward zero.


\vskip .1in

\noindent
{\bf
1. e4 e5
2. Nf3 Nf6
3. Nxe5 d6
4. Nd3?!
}

\vskip .1in
\noindent
While this is part of opening theory\footnote{ECO C42:
Petrov’s Defence, Karklins-Martinovsky Variation, to be precise.},
{\tt Stockfish} registered Carlsen's move as an inaccuracy,
preferring instead Nf3 or Nc4.

\vskip .2in
%\begin{figure}

\noindent
{\bf Position 1: After 4. Nd3 }

\noindent
{Biggest evaluation swing(s) (39cp)}

\noindent
Evaluation: +0.00
%\begin{center}

%\begin{small}
%\begin{tabular}{lcc}
%\toprule
%Metric & White & Black \\
%\midrule
%Space & 0.40 & 0.60 \\
%Mobility & 4.10 & 4.10 \\
%King Safety & 1.00 & 0.50 \\
%Threats & 0.80 & 2.80 \\
%\bottomrule
%\end{tabular}
%\end{small}

%\end{center}
%\end{figure}

\vskip .2in

\newpage

\subsection{The opening phase (5-15)}

\vskip .1in

\noindent
{\bf
4 ... Nxe4
5. Qe2 Qe7
6. Nf4 Nc6
7. Nd5 Nd4!
8. Nxe7! Nxe2!
9. Nd5! Nd4!
10. Na3 Ne6
11. f3! N4c5!
12. d4 Nd7
13. c3 c6!
14. Nf4! Nb6!
15. Bd3 d5
}
% fen = r1b1kb1r/pp3ppp/1np1n3/3p4/3P1N2/N1PB1P2/PP4PP/R1B1K2R w KQkq - 0 16

\vskip .1in
\noindent
This is all part of the opening. White is given a slight advantage by {\tt Stockfish}.

Figure \ref{fig:carlsen-caruana-position3}
tells us Caruana successfully seized territorial control (Space) in the opening up to move 15.
In Figure \ref{fig:carlsen-caruana-position4}, look at how the Mobility values drop
after the unusual early queen trade on move 8.


  \vskip .1in
\noindent
{\bf Position 2: After 15. ... d5 }
\begin{center}
  \chessboard[setfen=r1b1kb1r/pp3ppp/1np1n3/3p4/3P1N2/N1PB1P2/PP4PP/R1B1K2R w KQkq - 0 16]

  \vskip .1in
  After {\bf 15. ... d5}.
\end{center}


\subsection{The early middlegame (16-25)}

The game continued:
\vskip .2in

\noindent
{\bf
16. Nc2! Bd6
17. Nxe6 Bxe6
18. Kf2 h5
19. h4 Nc8
20. Ne3 Ne7
21. g3 c5!
22. Bc2 O-O
23. Rd1 Rfd8
24. Ng2 cxd4!
25. cxd4! Rac8
}

\vskip .1in
\noindent
This ends the early middlegame. The game is even, according to {\tt Stockfish}.

\vskip .1in
\noindent
{\bf Position 3: After 25. ... Rac8 }
\begin{center}
\begin{figure}[h!]
\begin{minipage}{\textwidth}
\begin{center}
%\begin{tabular}{c}
%\hspace{4.0 cm}
%\vspace{1.0 cm}
\chessboard[setfen=2rr2k1/pp2npp1/3bb3/3p3p/3P3P/5PP1/PPB2KN1/R1BR4 w - - 1 26]
\end{center}
%\\
\begin{center}
%\hspace{4.0 cm}
\includegraphics[height=7cm,width=12cm]{carlsen-caruana-game6-space-metric.png}
%\end{tabular}
\end{center}
\end{minipage}
\caption{(top) Position after 25 ... Rac8 in Carlsen-Caruana, Game 6 (2018). (bottom) The Space metric for that game. }
\label{fig:carlsen-caruana-position3}
\end{figure}
\end{center}

\newpage

\subsection{The later middlegame (26-33)}

The game continued:
\vskip .2in

\noindent
{\bf
26. Bb3! Nc6
27. Bf4! Na5
28. Rdc1! Bb4
29. Bd1 Nc4!
30. b3! Na3
31. Rxc8! Rxc8!
32. Rc1 Nb5
33. Rxc8+! Bxc8!
}

% fen = 2b3k1/pp3pp1/8/1n1p3p/1b1P1B1P/1P3PP1/P4KN1/3B4 w - - 0 34

\vskip .1in

\noindent
{\bf Position 4: After 33. ... Bxc8! }
\begin{center}
\begin{figure}[h!]
\begin{minipage}{\textwidth}
\begin{center}
%\begin{tabular}{c}
%\hspace{4.0 cm}
%\vspace{1.0 cm}
\chessboard[setfen=2b3k1/pp3pp1/8/1n1p3p/1b1P1B1P/1P3PP1/P4KN1/3B4 w - - 0 34]
\end{center}
%\\
\begin{center}
%\hspace{4.0 cm}
\includegraphics[height=7cm,width=12cm]{carlsen-caruana-game6-mobility-metric.png}
%\end{tabular}
\end{center}
\end{minipage}
\caption{(top) Position after 33 ... Bxc8 in Carlsen-Caruana, Game 6 (2018). (bottom) The Mobility metric for that game. }
\label{fig:carlsen-caruana-position4}
\end{figure}
\end{center}

\vskip .1in

\noindent
Note how the Mobility metric drops after these piece exchanges.

\newpage

\subsection{The early endgame (34-67)}

The game continued:
\vskip .2in

\noindent
{\bf
34. Ne3 Nc3!
35. Bc2 Ba3
36. Bb8 a6
37. f4 Bd7
38. f5 Bc6
39. Bd1! Bb2
40. Bxh5 Ne4+
41. Kg2 Bxd4!
42. Bf4 Bc5
43. Bf3! Nd2!
}

\vskip .1in

\noindent
This is a very interesting move, setting in motion a long combination,
dropping piece Mobility but escalating Threats.

\newpage

\noindent
{\bf Position 5: After 43. Bf3 }
\begin{center}
\begin{figure}[H]
\begin{minipage}{\textwidth}
\begin{center}
%\begin{tabular}{c}
%\hspace{4.0 cm}
%\vspace{1.0 cm}
\chessboard[setfen=6k1/1p3pp1/p1b5/2bp1P2/4nB1P/1P2NBP1/P5K1/8 b - - 3 43]
\end{center}
%\\
\begin{center}
%\hspace{4.0 cm}
\includegraphics[height=7cm,width=12cm]{carlsen-caruana-game6-threats-metric.png}
%\end{tabular}
\end{center}
\end{minipage}
\caption{(top) Position after 43 Bf3 and before 43 ... Nd2! in Carlsen-Caruana, Game 6 (2018). (bottom) The Threats metric for that game. }
\label{fig:carlsen-caruana-position5}
\end{figure}
\end{center}


\vskip .1in
\noindent
{\bf
44. Bxd5?
}

\vskip .1in

\noindent
{\tt Stockfish} didn't like Carlsen's move and suggest to consider instead Nd1 or Nf1 or Nc2.

\vskip .1in

\noindent
{\bf Position 6: After 44. Bxd5}

Evaluation: -0.81

%\chessboard[setfen=6k1/1p3pp1/p1b5/2bB1P2/5B1P/1P2N1P1/P2n2K1/8 b - - 0 44]

%\begin{small}
%\begin{tabular}{lcc}
%\toprule
%Metric & White & Black \\
%\midrule
%Space & 0.20 & 0.30 \\
%Mobility  & 2.70 & 2.30 \\
%King Safety & -0.80 & 0.20 \\
%Threats & 2.90 & 5.10 \\
%\bottomrule
%\end{tabular}
%\end{small}

\vskip .1in

\noindent
The game continued:
\vskip .2in

\noindent
{\bf
44...Bxe3!
45. Bxc6! Bxf4!
46. Bxb7! Bd6!
47. Bxa6 Ne4!
48. g4! Ba3!
49. Bc4 Kf8
50. g5 Nc3!
51. b4! Bxb4
52. Kf3 Na4!
53. Bb5 Nc5?!
54. a4?! f6
55. Kg4 Ne4
56. Kh5 Be1!
57. Bd3 Nd6
58. a5 Bxa5
59. gxf6! gxf6!
60. Kg6! Bd8
61. Kh7! Nf7
62. Bc4 Ne5!
63. Bd5 Ba5
64. h5! Bd2
65. Ba2 Nf3
66. Bd5 Nd4!
67. Kg6??
}

\vskip .1in
\noindent
{\tt Stockfish} considered this to be a losing mistake, suggesting instead: Be4.

\vskip .1in
\noindent
{\bf Position 7: After 67. Kg6}

Evaluation: -2.47

\chessboard[setfen=5k2/8/5pK1/3B1P1P/3n4/8/3b4/8 b - - 6 67]

%\begin{small}
%\begin{tabular}{lcc}
%\toprule
%Metric & White & Black \\
%\midrule
%Space & 0.00 & 0.10 \\
%Mobility (MG) & 1.60 & 1.80 \\
%King Safety & -2.20 & -1.60 \\
%Threats & 9.40 & 9.00 \\
%\bottomrule
%\end{tabular}
%\end{small}

\subsection{The later endgame (67-80)}

Black has a winning advantage at this point, according to
{\tt Stockfish}.

\vskip .1in

\noindent
{\bf
67...Bg5!
68. Bc4 Nf3??
}

\vskip .1in

\noindent
This is, according to {\tt Stockfish} a blunder, recommending instead Bh4.
According to {\tt Stockfish}, with depth=23, with this knight move Caruana missed a win.

\begin{center}
\begin{figure}[H]
\begin{minipage}{\textwidth}
\begin{center}
%\begin{tabular}{c}
%\hspace{4.0 cm}
%\vspace{1.0 cm}
\chessboard[setfen=5k2/8/5pK1/5PbP/2B5/5n2/8/8 w - - 9 69]
\end{center}
%\\
\begin{center}
%\hspace{4.0 cm}
\includegraphics[height=7cm,width=12cm]{carlsen-caruana-game6-king-safety-metric.png}
%\end{tabular}
\end{center}
\end{minipage}
\caption{(top) Position after 68 ... Nf3?? (moving from d4) in Carlsen-Caruana, Game 6 (2018).  (bottom) The King Safety metric for that game. }
\label{fig:carlsen-caruana-position8}
\end{figure}
\end{center}

\vskip .1in
After 18 Kf2, the White king enters the game early instead of castling, which
here created the opportunity for Black to make more threats.
Based on Figure \ref{fig:carlsen-caruana-position8} (bottom),
note the King Safety line drops into negative territory (-0.7 to -1.2) after the opening and stays there.
On the other hand, Black's King Safety was positive from moves 18 to 50 (after 49 ... Kf8).
Since Carlsen's king was more exposed to checks and needed protection, Caruana was able to
keep stronger positional evaluations, even with equal material.

\newpage

The game concluded this way:

\vskip .1in

\noindent
{\bf
69. Kh7! Ne5
70. Bb3 Ng4
71. Bc4 Ne3!
72. Bd3! Ng4!
73. Bc4! Nh6
74. Kg6! Ke7
75. Bb3! Kd6
76. Bc2! Ke5
77. Bd3 Kf4
78. Bc2! Ng4!
79. Bb3 Ne3
80. h6! Bxh6!
1/2-1/2
}

\subsection{FTI summary}

The various types of FTIs are designed to characterize the type of play.
Recall that the FTI metric is a linear combination of the four positional metrics,
Space (S), Mobility (M), King Safety (KS), and Threats (T):

\[
FTI = w_SS + w_M M + w_{KS}KS + w_T T.
\]
Think of the FTI as a coordinated pressure gauge.
Using it to predict a game's outcome requires a player to hold
positional dominance for 10 consecutive 'plies' (half-moves)
to predict a win. For the Carlsen-Caruana match, predicting the outcome
(win/loss/draw) using Fireteam Indices is a struggle because the game was so
well-played that the advantages according to most metrics were small.


In this case, CGA was run over all $12$ games Carlsen and Caruana played for the
2018 World Chess championships. The CGA program computes the
weights $(w_S, w_M, w_{KS}, w_T)$ such that, over all those games,
the FTI with those ``signature'' weights predicts the outcome
more successfully than any other choices. The outcome was

\[
(w_S, w_M, w_{KS}, w_T) = ( 0.0,  0.0,  0.0,  1.0 ),
\]
with $R^2 = 0$ (indicating a failure of linear regression to find a line
of ``best'' fit).

Either by looking at the plots for FTI1 and FTI5,
we see CGA suggests Caruana should have won the game and plays with
a solid, harmonious style.


\begin{center}
\begin{figure}[H]
\begin{minipage}{\textwidth}
\begin{center}
%\begin{tabular}{c}
%\hspace{4.0 cm}
%\vspace{1.0 cm}
\includegraphics[height=7cm,width=12cm]{carlsen-caruana-2018-FTI1.png}
\end{center}
%\\
\begin{center}
%\hspace{4.0 cm}
\includegraphics[height=7cm,width=12cm]{carlsen-caruana-2018-FTI2.png}
%\end{tabular}
\end{center}
\end{minipage}
\caption{Both plots are for game 6 of Carlsen-Caruana, 2018. (top) Fireteam index, FTI1 (harmonious). (bottom) Fireteam index, FTI2 (tactical). The second plot never exceeds the confidence threshold of $\pm 0.2$.}
\label{fig:carlsen-caruana-2018-FTI1+2}
\end{figure}
\end{center}

\begin{center}
\begin{figure}[H]
\begin{minipage}{\textwidth}
\begin{center}
%\begin{tabular}{c}
%\hspace{4.0 cm}
%\vspace{1.0 cm}
\includegraphics[height=7cm,width=10cm]{carlsen-caruana-2018-FTI3.png}
\end{center}
%\\
\begin{center}
%\hspace{4.0 cm}
\includegraphics[height=7cm,width=12cm]{carlsen-caruana-2018-FTI4.png}
%\end{tabular}
\end{center}
%\\
\begin{center}
%\hspace{4.0 cm}
\includegraphics[height=7cm,width=12cm]{carlsen-caruana-2018-FTI5.png}
%\end{tabular}
\end{center}
\end{minipage}
\caption{These plots are also for game 6 of Carlsen-Caruana, 2018. (top) Fireteam index, FTI3 (strategic). (middle) Fireteam index, FTI4 (dynamic). (bottom) Fireteam index, FTI5 (solid).}
\label{fig:carlsen-caruana-2018-FTI3+4+5}
\end{figure}
\end{center}

\vskip .1in
Of all the FTI plots in Figures \ref{fig:carlsen-caruana-2018-FTI1+2}
and \ref{fig:carlsen-caruana-2018-FTI3+4+5}, the first and last one 
most strongly suggests a win for Black. The other FTIs also suggest a
Black win, but with less confidence.

% =============================================================================
% Chess Game Analyzer - README Documentation
% Suitable for inclusion as an appendix in a manual
% =============================================================================

\chapter{Appendix: Chess Game Analyzer module reference}
\label{appendix:cga-readme}

This is the readme file for the CGA program \cite{CGA26} generated by {\tt claude} (Anthropic),
with some editing by me.

\section{Overview}

The Chess Game Analyzer is a Python module for deep positional analysis of chess games using Stockfish, with comprehensive \LaTeX{} report generation and win prediction metrics.

Chess Game Analyzer goes beyond simple engine evaluation to provide interpretable positional metrics for every move of a chess game. It computes Space control, Mobility, King Safety, and Threats directly from board positions, then combines these into \textbf{Fireteam Index (FTI)} metrics that can predict game outcomes based on sustained positional advantages.

The module generates professional \LaTeX{} reports with annotated games, critical position diagrams, matplotlib plots, and detailed methodology explanations---suitable for chess study, coaching, or research.


% -----------------------------------------------------------------------------
\section{Background for Newcomers}
% -----------------------------------------------------------------------------

\subsection{What is Stockfish?}

{\tt Stockfish}\footnote{\url{https://stockfishchess.org/}} is the world's strongest open-source chess engine. It analyzes positions and calculates the best moves by searching millions of positions per second. This module uses Stockfish to evaluate positions, but computes the positional breakdown (Space, Mobility, etc.)\ directly using the {\tt python-chess} library.

\textbf{Stockfish must be installed on your system.} Download it from the
official website and note the installation path.

\subsection{What is a Centipawn?}

Chess engines measure advantages in \textbf{centipawns (cp)}---hundredths of a pawn. An evaluation of $+100$ means White is ahead by roughly one pawn's worth of material or positional compensation. An evaluation of $-250$ means Black has about a 2.5 pawn advantage.

\begin{center}
\begin{tabular}{ll}
\toprule
\textbf{Evaluation} & \textbf{Meaning} \\
\midrule
$0$ to $\pm 50$ & Roughly equal \\
$\pm 50$ to $\pm 150$ & Slight advantage \\
$\pm 150$ to $\pm 300$ & Clear advantage \\
$\pm 300$ to $\pm 500$ & Winning advantage \\
$\pm 500+$ & Decisive advantage \\
$\pm 10000$ & Checkmate is forced \\
\bottomrule
\end{tabular}
\end{center}

\subsection{Python Requirements}

The module requires:
\begin{itemize}
    \item \textbf{Python 3.8 or higher}
    \item \textbf{python-chess}: The chess library that handles board representation, move generation, and PGN parsing
    \item \textbf{matplotlib} (optional): For generating plots; the module works without it using ASCII plots
\end{itemize}

To install the dependencies:
\begin{verbatim}
pip install python-chess matplotlib
\end{verbatim}

% -----------------------------------------------------------------------------
\section{Installation}
% -----------------------------------------------------------------------------

\begin{enumerate}
    \item \textbf{Install Stockfish} from \url{https://stockfishchess.org/download/}
    
    \item \textbf{Clone the repository} or download \texttt{chess\_game\_analyzer.py}
    
    \item \textbf{Install Python dependencies}:
\begin{verbatim}
pip install python-chess matplotlib
\end{verbatim}
    
    \item \textbf{Verify Stockfish path}:
\begin{verbatim}
# Linux/Mac - find where Stockfish is installed
which stockfish

# Or test directly
stockfish
\end{verbatim}
\end{enumerate}

% -----------------------------------------------------------------------------
\section{Quick Start}
% -----------------------------------------------------------------------------

\subsection{Command Line}

\begin{verbatim}
# Analyze a single game
python chess_game_analyzer.py game.pgn -o analysis.tex

# Analyze with custom Stockfish path and deeper analysis
python chess_game_analyzer.py game.pgn -o analysis.tex \
    --stockfish /usr/local/bin/stockfish \
    --depth 22 \
    --time 2.0

# Analyze multiple games as a book
python chess_game_analyzer.py tournament.pgn -o book.tex --book \
    --book-title "My Tournament Games" \
    --book-author "Your Name"
\end{verbatim}

\subsection{Python API}

\begin{verbatim}
from chess_game_analyzer import (
    EnhancedGameAnalyzer,
    EnhancedLaTeXReportGenerator,
    classify_game_character
)

# Analyze a game
with EnhancedGameAnalyzer("/usr/local/bin/stockfish", 
                          depth=20, time_limit=1.0) as analyzer:
    result = analyzer.analyze_game("game.pgn")

# Access move-by-move data
for move in result.moves:
    pos = move.positional_eval
    print(f"Move {move.ply}: {move.move_san}")
    print(f"  Space: W={pos.space_white:.2f}, B={pos.space_black:.2f}")
    print(f"  Mobility: W={pos.mobility_white:.2f}, B={pos.mobility_black:.2f}")
    print(f"  Eval: {move.eval_after/100:+.2f} pawns")

# Generate LaTeX report
latex = EnhancedLaTeXReportGenerator.generate_report(
    result,
    include_plots=True,
    include_methodology=True
)
with open("analysis.tex", "w") as f:
    f.write(latex)
\end{verbatim}

% -----------------------------------------------------------------------------
\section{Positional Metrics Explained}
% -----------------------------------------------------------------------------

\subsection{Space}

\textbf{Space} measures territorial control, particularly in the center and the opponent's half of the board.

The computation counts squares on central files (c--f) in ranks 2--4 (for White) or 5--7 (for Black) that are:
\begin{enumerate}
    \item Attacked by at least one friendly pawn, {\it and}
    \item Not occupied by an enemy piece
\end{enumerate}

The raw count is divided by 10 to normalize the metric. Values typically range from 0.1 to 0.7.

\paragraph{Example from Estrin vs Berliner (1965), move 20 (10...e4):}
\begin{itemize}
    \item White Space: 0.30
    \item Black Space: 0.30
\end{itemize}
At this point in a sharp Two Knights Defense, space is roughly equal, but the position is about to explode tactically.

\subsection{Mobility}

\textbf{Mobility} counts the number of safe squares available to each piece, weighted by piece type.

``Safe'' squares exclude:
\begin{itemize}
    \item Squares attacked by enemy pawns
    \item Squares that would result in material loss through exchange
\end{itemize}

Mobility in enemy territory is weighted higher. The metric captures piece activity and coordination.

\paragraph{Example from Estrin vs Berliner, move 30 (15...Bxf3):}
\begin{itemize}
    \item White Mobility: 1.95
    \item Black Mobility: 2.5
\end{itemize}
After White's blunder 15.Be2??, Black's pieces have become significantly more active.

\subsection{King Safety}

\textbf{King Safety} combines several factors:

\begin{enumerate}
    \item \textbf{Pawn shield strength}: Pawns on the 2nd/3rd rank in front of the castled king
    \item \textbf{King geometry}: Distance of enemy pieces to your king
    \item \textbf{Attack units}: Enemy pieces attacking squares near your king
\end{enumerate}

Positive values indicate a safe king; negative values indicate danger.

\begin{example} From Estrin vs Berliner, move 52 (26...hxg2+):
\begin{itemize}
    \item White King Safety: $-1.40$
    \item Black King Safety: $-0.70$
\end{itemize}
White's king is significantly more exposed after the pawn storm breaks through.
\end{example}

\subsection{Threats}

\textbf{Threats} evaluates the tactical tension in the position:

\begin{itemize}
    \item \textbf{Hanging pieces}: Pieces attacked but not defended (weighted by piece value $\times 2$)
    \item \textbf{Attacks by lower-value pieces}: e.g., a knight attacking a queen
    \item \textbf{King zone pressure}: Attacks on squares near the enemy king
    \item \textbf{Safe checks}: Available checking moves that don't lose material
    \item \textbf{Weak squares (holes)}: Squares that can never be defended by pawns
\end{itemize}

Raw Threats values typically range from 1 to 30. In the Fireteam Index calculations, Threats is scaled down (divided by 20) to be comparable with the other metrics.

\begin{example} From Estrin vs Berliner, move 24.Nc4: The Threats graph
  has spikes and valleys in this part of the game.
\begin{itemize}
    \item White Threats: 21.80
    \item Black Threats: 8.50
\end{itemize}
White has significant pressure, but Black's position is resilient.
\end{example}

% -----------------------------------------------------------------------------
\section{The Fireteam Index (FTI)}
% -----------------------------------------------------------------------------

\subsection{Background: Chess Metrics and Win Prediction}

Traditional chess analysis focuses on \textbf{centipawn evaluation}---a single number indicating who stands better. But this doesn't tell you \emph{why} someone is winning or \emph{how} the advantage manifests.

The \textbf{Fireteam Index} is inspired by Rithmomachia,\footnote{Rithmomachia, also known as the ``philosopher's game,'' was a medieval mathematical board game popular in European monasteries and universities from the 11th to 16th centuries.} a medieval mathematical board game where one victory condition is establishing a ``well-coordinated fireteam'' of pieces deep in enemy territory forming arithmetic, geometric, or harmonic progressions.

Similarly, in chess, sustained positional dominance across multiple factors---space, activity, king safety, and threats---often precedes victory. The FTI combines these factors into a single metric that tracks this multi-dimensional advantage over time.

\subsection{FTI Signatures}

Different players have different styles. A tactical player might win through direct threats, while a positional player grinds down opponents through space and piece activity. The module provides five FTI ``signatures'' with different weightings:

\begin{center}
\begin{tabular}{lccccc}
\toprule
\textbf{FTI Variant} & \textbf{Space} & \textbf{Mobility} & \textbf{King Safety} & \textbf{Threats} & \textbf{Style} \\
\midrule
FTI$_1$ (Harmonious) & 0.25 & 0.25 & 0.25 & 0.25 & Balanced \\
FTI$_2$ (Tactical) & 0.60 & 0.10 & 0.00 & 0.30 & Space + threats \\
FTI$_3$ (Strategic) & 0.70 & 0.10 & 0.10 & 0.10 & Spatial domination \\
FTI$_4$ (Dynamic) & 0.20 & 0.60 & 0.00 & 0.20 & Piece activity \\
FTI$_5$ (Solid) & 0.30 & 0.20 & 0.50 & 0.00 & Defensive play \\
\bottomrule
\end{tabular}
\end{center}

The FTI formula for a given player is:
\begin{equation}
\text{FTI} = w_S \cdot \Delta S + w_M \cdot \Delta M + w_K \cdot \Delta K + w_T \cdot \Delta T'
\end{equation}
where each $\Delta$ is (Player's value $-$ Opponent's value) and $T' = T / 20$ (scaled for comparability).

\subsection{Prediction Algorithms}

Three algorithms use the FTI to predict game outcomes:

\paragraph{1. Per-Ply (Raw).}
\begin{itemize}
    \item Examines raw FTI at each ply after move 8 (ply 16)
    \item Predicts {\it win} if FTI $> \epsilon$ (default 0.2) for 10 consecutive plies
    \item Captures sudden, decisive advantages
\end{itemize}

\paragraph{2. Windowed (Smoothed).}
\begin{itemize}
    \item Computes a 10-ply rolling average of FTI
    \item Applies the same streak detection to smoothed values
    \item Filters out tactical noise; captures sustained trends
\end{itemize}

\paragraph{3. Balance (Confidence-based).}
\begin{itemize}
    \item Counts plies where FTI $> \epsilon$ versus FTI $< -\epsilon$
    \item Measures volatility (sign changes)
    \item Computes confidence $= |B_n| \times (1 - V/V_{\rm max})$
    \item Predicts {\it win/loss} if confidence $> 0.15$; otherwise {\it draw}.
    \item Best at detecting drawn games and balanced struggles.
\end{itemize}

% -----------------------------------------------------------------------------
\section{Concrete Example: Estrin vs Berliner (1965)}
% -----------------------------------------------------------------------------

Let us walk through the analysis of a famous game from the 5th World Correspondence Chess Championship---\textbf{Yakov Estrin vs Hans Berliner}, featuring the sharp Two Knights Defense.

\subsection{The Game (C57 -- Two Knights Defense)}

As we know form the previous chapter, the first few moves of the game are
as follows.

\begin{quote}
{\bf 1.e4 e5 2.Nf3 Nc6 3.Bc4 Nf6 4.Ng5 d5 5.exd5 b5!?}
\end{quote}

Berliner plays the Ulvestad Variation (5...b5), a sharp gambit that had been analyzed deeply in correspondence play.

\subsection{Analysis Output}

\paragraph{Player Statistics.}

\begin{center}
\begin{tabular}{lcc}
\toprule
\textbf{Metric} & \textbf{Estrin (White)} & \textbf{Berliner (Black)} \\
\midrule
Accuracy & 87.5\% & 92.1\% \\
Avg centipawn loss & 26.6 & 16.5 \\
Best moves & 18 & 21 \\
Blunders & 2 & 0 \\
\bottomrule
\end{tabular}
\end{center}

\paragraph{Game Character Classification.}
\begin{itemize}
    \item Evaluation range: $m_1 = -7.04$, $m_2 = +2.24$
    \item Spread: $d = 9.28$ pawns
    \item Classification: \textbf{Chaotic seesaw battle}
\end{itemize}

The game swung wildly---White had $+2.24$ at one point, but Black eventually achieved $-7.04$.

\subsection{Positional Metrics Summary}

\begin{center}
\begin{tabular}{lcc}
\toprule
\textbf{Metric} & \textbf{White Average} & \textbf{Black Average} \\
\midrule
Space & 0.30 & 0.22 \\
Mobility & 1.51 & 1.67 \\
King Safety & $-0.25$ & $-0.22$ \\
Threats & 6.35 & 5.91 \\
\bottomrule
\end{tabular}
\end{center}

Despite the sharp tactical nature, Berliner maintained slightly better piece activity throughout.


\subsection{Move-by-Move Data Extract}

The module stores raw positional data in a machine-readable format:

\begin{verbatim}
% ply, SAN, eval_cp, space_w, space_b, mob_w, mob_b, ks_w, ks_b, 
%      threats_w, threats_b
% 29, Be2, -126, 0.300, 0.300, 5.000, 4.600, 0.000, 0.100, 3.800, 4.700
% 30, Bxf3, -204, 0.200, 0.300, 3.900, 5.000, -0.200, 0.500, 7.600, 6.000
% 31, Bxf3, -237, 0.300, 0.300, 4.000, 4.400, 0.200, 0.500, 7.400, 5.200
\end{verbatim}

This data can be parsed for further computational analysis, machine learning, or statistical studies.

% -----------------------------------------------------------------------------
\section{Command Line Interface}
% -----------------------------------------------------------------------------

\begin{verbatim}
python chess_game_analyzer.py [PGN_FILE] [OPTIONS]
\end{verbatim}

\subsection{Core Options}

\begin{center}
\begin{tabular}{lll}
\toprule
\textbf{Option} & \textbf{Description} & \textbf{Default} \\
\midrule
\texttt{-o, --output FILE} & Output \LaTeX{} file path & stdout \\
\texttt{--json-output FILE} & Output raw analysis as JSON & none \\
\texttt{-s, --stockfish PATH} & Path to Stockfish executable & \texttt{/usr/games/stockfish} \\
\texttt{-d, --depth N} & Analysis depth (higher = slower) & 20 \\
\texttt{-t, --time SECS} & Time limit per position & 1.0 \\
\texttt{-q, --quiet} & Suppress progress messages & false \\
\bottomrule
\end{tabular}
\end{center}

\subsection{Report Options}

\begin{center}
\begin{tabular}{ll}
\toprule
\textbf{Option} & \textbf{Description} \\
\midrule
\texttt{--no-diagrams} & Don't include chess board diagrams \\
\texttt{--no-methodology} & Don't include methodology explanation section \\
\texttt{--no-plots} & Don't include matplotlib plots \\
\texttt{--ascii-plots} & Include ASCII plots (works without matplotlib) \\
\texttt{--plot-dir DIR} & Directory for plot output files \\
\bottomrule
\end{tabular}
\end{center}

\subsection{Fireteam Index Prediction (Single Game)}

\begin{center}
\begin{tabular}{ll}
\toprule
\textbf{Option} & \textbf{Description} \\
\midrule
\texttt{--prediction W|B} & Include FTI prediction for White (W) or Black (B) \\
\texttt{--prediction-name NAME} & Player name for display (e.g., ``Berliner'') \\
\bottomrule
\end{tabular}
\end{center}

\subsection{Book Mode (Multiple Games)}

\begin{center}
\begin{tabular}{ll}
\toprule
\textbf{Option} & \textbf{Description} \\
\midrule
\texttt{--book} & Analyze all games in PGN as a book \\
\texttt{--book-title TITLE} & Title for the book \\
\texttt{--book-author AUTHOR} & Author name \\
\texttt{--prediction-winner} & Track winner in decisive games (default: true) \\
\texttt{--prediction-name NAME} & Track specific player by name \\
\texttt{--no-prediction} & Disable FTI prediction section entirely \\
\bottomrule
\end{tabular}
\end{center}

\subsection{Examples}

\begin{verbatim}
# Deep analysis with plots
python chess_game_analyzer.py game.pgn -o report.tex \
    --stockfish /opt/stockfish/stockfish \
    --depth 24 --time 3.0

# Quick analysis without graphics
python chess_game_analyzer.py game.pgn -o report.tex \
    --depth 16 --time 0.5 --no-plots --ascii-plots

# Analyze Berliner's games as a collection
python chess_game_analyzer.py berliner_games.pgn -o berliner_book.tex \
    --book \
    --book-title "Berliner's World Championship Games" \
    --book-author "Analysis by Stockfish 17" \
    --prediction-name "Berliner"

# JSON output for programmatic access
python chess_game_analyzer.py game.pgn --json-output analysis.json
\end{verbatim}

% -----------------------------------------------------------------------------
\section{Python API}
% -----------------------------------------------------------------------------

\subsection{Main Classes}

\paragraph{EnhancedGameAnalyzer.}
The core analysis engine. Use as a context manager to ensure proper cleanup.

\begin{verbatim}
with EnhancedGameAnalyzer(
    stockfish_path="/usr/local/bin/stockfish",
    depth=20,           # Search depth (plies)
    time_limit=1.0,     # Seconds per position
    extract_positional=True  # Compute Space/Mobility/etc.
) as analyzer:
    result = analyzer.analyze_game("game.pgn")
    # or analyze_all_games() for multi-game PGNs
\end{verbatim}

\paragraph{EnhancedGameAnalysisResult.}
Contains all analysis data:
\begin{itemize}
    \item \texttt{moves}: List of \texttt{EnhancedMoveAnalysis} objects
    \item \texttt{white\_stats}, \texttt{black\_stats}: Accuracy and move classification counts
    \item \texttt{brilliant\_sacrifices}: Detected sacrifices
    \item \texttt{critical\_positions}: Key moments for diagrams
    \item \texttt{game\_character}: Classification (balanced/tactical/chaotic, etc.)
    \item \texttt{positional\_summary}: Aggregated Space/Mobility/etc.\ statistics
\end{itemize}

\paragraph{PositionalEvaluation.}
Attached to each move, containing:
\begin{itemize}
    \item \texttt{space\_white}, \texttt{space\_black}
    \item \texttt{mobility\_white}, \texttt{mobility\_black}
    \item \texttt{king\_safety\_white}, \texttt{king\_safety\_black}
    \item \texttt{threats\_white}, \texttt{threats\_black}
\end{itemize}

\subsection{Utility Functions}

\begin{verbatim}
from chess_game_analyzer import (
    parse_raw_positional_data,
    compute_fireteam_index,
    predict_outcome_per_ply,
    predict_outcome_windowed,
    predict_outcome_balance,
    classify_game_character
)

# Parse raw data from LaTeX file
data = parse_raw_positional_data(tex_content, game_number=7)

# Compute FTI from positional data
fti_data = compute_fireteam_index(data)

# Run predictions with custom weights
from chess_game_analyzer import FTI2_WEIGHTS  # Tactical signature

result = predict_outcome_per_ply(
    data,
    player_color="B",
    player_name="Berliner",
    weights=FTI2_WEIGHTS,
    cutoff_ply=16,
    streak_length=10,
    epsilon=0.2
)
print(f"Prediction: {result.prediction}")
print(f"Max streak: {result.max_streak} plies")
\end{verbatim}

% -----------------------------------------------------------------------------
\section{Parameter Reference}
% -----------------------------------------------------------------------------

\subsection{Analysis Parameters}

\begin{center}
\begin{tabular}{llll}
\toprule
\textbf{Parameter} & \textbf{Type} & \textbf{Default} & \textbf{Description} \\
\midrule
\texttt{stockfish\_path} & str & \texttt{/usr/games/stockfish} & Path to Stockfish binary \\
\texttt{depth} & int & 20 & Search depth in plies \\
\texttt{time\_limit} & float & 1.0 & Maximum seconds per position \\
\texttt{min\_diagram\_spacing} & int & 6 & Min ply distance between diagrams \\
\texttt{top\_n\_swings} & int & 2 & ``Biggest swing'' positions to include \\
\bottomrule
\end{tabular}
\end{center}

\subsection{FTI Prediction Parameters}

\begin{center}
\begin{tabular}{llll}
\toprule
\textbf{Parameter} & \textbf{Type} & \textbf{Default} & \textbf{Description} \\
\midrule
\texttt{cutoff\_ply} & int & 16 & Opening plies to skip \\
\texttt{streak\_length} & int & 10 & Consecutive plies for WIN \\
\texttt{epsilon} & float & 0.2 & Min FTI for ``meaningful advantage'' \\
\texttt{window\_size} & int & 10 & Rolling average window \\
\texttt{weights} & tuple & $(0.25, 0.25, 0.25, 0.25)$ & $(w_S, w_M, w_K, w_T)$ \\
\texttt{v\_ratio} & float & 0.20 & Volatility threshold ratio \\
\texttt{confidence\_threshold} & float & 0.15 & Min confidence for WIN/LOSS \\
\bottomrule
\end{tabular}
\end{center}

\subsection{Available Weight Presets}

\begin{verbatim}
from chess_game_analyzer import (
    FTI1_WEIGHTS,  # (0.25, 0.25, 0.25, 0.25) - Harmonious
    FTI2_WEIGHTS,  # (0.60, 0.10, 0.00, 0.30) - Tactical
    FTI3_WEIGHTS,  # (0.70, 0.10, 0.10, 0.10) - Strategic
    FTI4_WEIGHTS,  # (0.20, 0.60, 0.00, 0.20) - Dynamic
    FTI5_WEIGHTS,  # (0.30, 0.20, 0.50, 0.00) - Solid
)
\end{verbatim}

\subsection{Move Classification Thresholds}

\begin{center}
\begin{tabular}{ll}
\toprule
\textbf{Classification} & \textbf{Centipawn Loss} \\
\midrule
Best & 0 (matches engine's first choice) \\
Excellent & 0--10 cp \\
Good & 10--25 cp \\
Inaccuracy & 25--50 cp \\
Mistake & 50--150 cp \\
Blunder & 150+ cp \\
\bottomrule
\end{tabular}
\end{center}

% -----------------------------------------------------------------------------
\section{Output Formats}
% -----------------------------------------------------------------------------

\subsection{\LaTeX{} Report}

The primary output is a \LaTeX{} document that can be compiled with \texttt{pdflatex} or \texttt{xelatex}. It includes:

\begin{enumerate}
    \item \textbf{Game Information}: Event, players, date, opening
    \item \textbf{Player Statistics}: Accuracy, move counts by classification
    \item \textbf{Positional Analysis}: Space/Mobility/KingSafety/Threats summaries with plots
    \item \textbf{Game Character Classification}: Balanced/tactical/chaotic assessment
    \item \textbf{Annotated Game}: Move-by-move with NAG symbols and comments
    \item \textbf{Critical Positions}: Diagrams at key moments with best alternatives
    \item \textbf{Methodology Section}: (optional) Detailed explanation of all metrics
    \item \textbf{Fireteam Index Prediction}: (optional) Win prediction analysis with plots
\end{enumerate}

Required \LaTeX{} packages: \texttt{xskak}, \texttt{chessboard}, \texttt{amsmath}, \texttt{pgfplots}, \texttt{graphicx}, \texttt{hyperref}.

\subsection{JSON Output}

Raw analysis data in JSON format for programmatic access:

\begin{verbatim}
{
  "white": "Estrin",
  "black": "Berliner",
  "result": "0-1",
  "moves": [
    {
      "ply": 1,
      "move_san": "e4",
      "eval_after": 15,
      "classification": "best",
      "positional_eval": {
        "space_white": 0.40,
        "space_black": 0.40,
        "mobility_white_mg": 3.90,
        ...
      }
    },
    ...
  ],
  "white_stats": {"accuracy": 87.5, ...},
  "black_stats": {"accuracy": 92.1, ...}
}
\end{verbatim}

\subsection{Raw Positional Data}

Appended after \verb|\end{document}| in \LaTeX{} files, machine-readable CSV format:

\begin{verbatim}
% GAME_DATA_START
% 1, e4, 15, 0.400, 0.400, 3.900, 3.300, 1.000, 1.500, 0.800, 0.800
% 2, e5, 24, 0.400, 0.400, 3.900, 3.900, 1.000, 1.000, 0.800, 0.800
...
% GAME_DATA_END
\end{verbatim}

Fields: ply, SAN, eval\_cp, space\_w, space\_b, mob\_w, mob\_b, ks\_w, ks\_b, threats\_w, threats\_b.

% -----------------------------------------------------------------------------
\section{License}
% -----------------------------------------------------------------------------

This software is released under your choice of either the \textbf{MIT License} or the \textbf{Modified BSD License}.

\begin{quote}
Copyright \textcopyright{} 2026, David Joyner
\end{quote}

% -----------------------------------------------------------------------------
\section{Acknowledgments}
% -----------------------------------------------------------------------------

\begin{itemize}
    \item \textbf{Stockfish} (\url{https://stockfishchess.org/}) --- The chess engine powering the evaluations
    \item \textbf{python-chess} (\url{https://python-chess.readthedocs.io/}) --- The excellent chess library
    \item The concept of the Fireteam Index is inspired by \textbf{Rithmomachia}, the medieval ``philosopher's game''
\end{itemize}

\chapter{References}

\begin{thebibliography}{999}

\bibitem[Du17]{Du17} A. Dunne, {\it Hans Berliner, 1929-2017},
Chess Life, May 2017, pp43-45.

\bibitem[Be65]{Be65} {\it Quantitative-Analytic Reports on Chess:
  Berliner’s games, 5th World Ch. ICCF 1965-1968}
  ({\small Generated by Stockfish 17.1 on January 28, 2026}),
  the QARC pdf of Berliner's games in the
  5th World Correspondence Chess Championship, ICCF, 1965.
Available at:
\newline
\url{https://github.com/wdjoyner/qarc}

\bibitem[Be99]{Be99} H. Berliner, {\bf The system}, Gambit Publ., 1999.

\bibitem[Be05]{Be05} Interview with Gardner Hendrie on 2005-03-07,
{\it Oral history of Hans Berliner},  
CHM Reference number X3131.2005, Computer History Museum, 2005.  


\bibitem[CC18]{CC18} {\it Quantitative-Analytic Reports on Chess:
   Carlsen vs Caruana, 2018 World Championship, London}
  ({\small Generated by Stockfish 17.1 on January 29, 2026}),
  the QARC pdf of all 12 games in the
  2018 World Chess Championship, FIDE.
Available at:
\newline
\url{https://github.com/wdjoyner/qarc}

\bibitem[CGA26]{CGA26} D. Joyner, the Chess Game Analyzer, version 6y.
Available at:
\newline
\url{https://github.com/wdjoyner/chess}


\bibitem[FEN]{FEN} D. Forsyth and S. Edwards,
  {\it Portable Game Notation Specification and Implementation Guide},
Available at:
\newline
\url{https://github.com/fsmosca/PGN-Standard/blob/master/PGN-Standard.txt}
  
\bibitem[PV]{PV} {\it Chess piece relative value}, wikipedia.
Available at:
\newline
\url{https://en.wikipedia.org/wiki/Chess_piece_relative_value}
  
\bibitem[Ri25]{Ri25}
David Joyner,
{\it Rithmomachia: is recreational arithmetic}, monograph, 2025. 
Available at {\tt amazon.com} or in the docs directory at:
\newline
\url{https://github.com/wdjoyner/rithmomachia}

\bibitem[To12]{To12} J. Togyer,  {\it The Iconoclast}, in
{\bf The Link: The Magazine of the Carnegie Mellon University School of Computer Science}, May 7, 2012.

\end{thebibliography}
\end{document}


